\documentclass[11pt]{article}
\usepackage{amsmath, amssymb, geometry}
\geometry{margin=1in}

\begin{document}

\begin{center}
    {\Large \textbf{Deriving the Optimal Tariff Formula}} \\[0.3cm]
\end{center}

\section*{1. Setup}

Consider a large importing country. Let
\[
D(p) \quad \text{be domestic demand}, \qquad 
S(p) \quad \text{be domestic supply},
\]
and
\[
m(p)=D(p)-S(p)
\]
be imports. Domestic consumers pay the price $p$, while foreign exporters receive
\[
p_W = p_W(m),
\]
where $p_W'(m)>0$, reflecting an upward-sloping foreign export supply.

The government imposes an ad-valorem tariff such that
\[
t=\frac{p-p_W}{p_W}.
\]

\section*{2. Consumer Surplus, Producer Surplus, and Tariff Revenue}

Consumer surplus is the area under demand and above the price:
\[
CS(p)
= \int_0^{D(p)} P_D(q)\, dq - pD(p),
\]
where $P_D(D(p))=p$ is inverse demand.

Producer surplus is the area above supply and below the price:
\[
PS(p)
= pS(p) - \int_0^{S(p)} P_S(q)\, dq,
\]
where $P_S(S(p))=p$ is inverse supply.

Tariff revenue is
\[
R(p)
= (p-p_W)\,m(p)
= (p-p_W)\,[D(p)-S(p)].
\]

Total national welfare is
\[
W(p)=CS(p)+PS(p)+R(p).
\]

\section*{3. Combine CS, PS, and R}

Add the three expressions:
\begin{align*}
W(p) &= \left[\int_0^{D} P_D(q)\,dq - pD\right]
     + \left[pS - \int_0^{S} P_S(q)\,dq\right]
     + (p-p_W)(D-S).
\end{align*}

Group terms:
\[
W(p)
= \int_0^{D}P_D(q)\,dq 
- \int_0^{S}P_S(q)\,dq
+ [-pD + pS + (p-p_W)(D-S)].
\]

Expand the price bracket:
\begin{align*}
-pD + pS + (p-p_W)(D-S)
&= -pD + pS + pD - pS - p_W(D-S) \\
&= -p_W(D-S).
\end{align*}

Thus:
\[
\boxed{
W(p)
= \int_{0}^{D(p)} P_D(q)\,dq
- \int_{0}^{S(p)} P_S(q)\,dq
- p_W(m(p))\, m(p).
}
\]

\section*{4. Differentiating Welfare}

Apply Leibniz's Rule.

\subsection*{4.1 Derivative of the consumer surplus integral}
\[
\frac{d}{dp}\int_{0}^{D(p)}P_D(q)\,dq
= P_D(D(p))D'(p) = pD'(p).
\]

\subsection*{4.2 Derivative of the producer surplus integral}
\[
\frac{d}{dp}\left[-\int_{0}^{S(p)}P_S(q)\,dq\right]
= -P_S(S(p))S'(p)= -pS'(p).
\]

\subsection*{4.3 Derivative of the import-cost term}

Let $m(p)=D(p)-S(p)$ and $m'(p)=D'(p)-S'(p)$. Then
\[
\frac{d}{dp}[p_W(m)m]
= \frac{dp_W}{dm}m'(p)m(p) + p_W(m)m'(p).
\]

Thus,
\[
\frac{d}{dp}[-p_W(m)m]
= -\frac{dp_W}{dm}m'(p)m(p) - p_W(m)m'(p).
\]

\section*{5. Total derivative}

Add the three components:
\begin{align*}
\frac{dW}{dp}
&= pD'(p) - pS'(p)
   - m'(p)\frac{dp_W}{dm}m(p)
   - p_W(m)m'(p) \\
&= p m'(p) - p_W(m)m'(p)
 - m(p)\frac{dp_W}{dm}m'(p).
\end{align*}

Factor out $m'(p)$:
\[
\boxed{
\frac{dW}{dp}
= m'(p)\left[p - p_W(m) - m(p)\frac{dp_W}{dm}\right].
}
\]

The first-order condition (for $m'(p)\neq 0$) is:
\[
p - p_W(m) - m\frac{dp_W}{dm} = 0.
\]

Thus the optimal tariff condition is:
\[
\boxed{
p = p_W(m) + m\,\frac{dp_W}{dm}.
}
\]

\section*{6. Elasticity Form}

Let the foreign export supply elasticity be
\[
\eta
= \frac{d\ln m}{d\ln p_W}
= \frac{p_W}{m}\left(\frac{dp_W}{dm}\right)^{-1}.
\]

Rearrange the optimality condition:
\[
p - p_W = m\frac{dp_W}{dm}
= \frac{p_W}{\eta}.
\]

Therefore:
\[
\boxed{
\frac{p-p_W}{p_W} = \frac{1}{\eta}.
}
\]

This is the \textbf{optimal tariff formula}.

\end{document}
